\documentclass[11pt,a4paper]{article}
\usepackage{turkos}

\begin{document}
\phaseheader{8}{Processes, Context Switching and Scheduling}{Run many things at once on one CPU: kernel threads, the context switch, and a preemptive scheduler driven by the timer}{48}{56}

\ataglance{3 weeks}{5}{Phase 7: \code{kmalloc}, \code{list.h}, \code{make test}; Phase 4: the timer interrupt}{Kernel threads, each
with its own stack and guard page, a context switch in about ten lines of assembly, cooperative \code{yield}, then
preemptive round-robin scheduling on the timer, \code{sleep_ms} with proper blocking, an idle task that halts the CPU,
task exit with reaping, and a \code{ps} command. Optionally, a multi-level feedback queue.}

\section{Why this phase}

This is where Turk-OS stops being a program and becomes an operating system. A single CPU can only run one
instruction stream at a time, but by switching between streams quickly enough, the kernel creates the illusion that
many run at once. \OSTEP calls this \emph{virtualizing the CPU}, and it is the first of its three big ideas.

You will build it in two layers, as the books do. The \textbf{mechanism} is the context switch: saving everything
that describes one running thread and restoring another. The \textbf{policy} is the scheduler: deciding which thread
runs next and for how long. You start with kernel threads, which all share the kernel's address space; in Phase 10
the same machinery switches between user processes with separate address spaces. Take your time with the context
switch. It is short, it is subtle, and every later phase depends on it.

\section{Concepts you need}

\subsection{The process abstraction}

\OSTEP chapter 4 defines a process by its \textbf{machine state}: the registers (including the instruction and stack
pointers), the memory it can address, and the I/O resources it holds. The kernel keeps this state in a
\textbf{process control block} (PCB), which Turk-OS calls \code{struct task}, and tracks every task's
\textbf{state} (Figure~\ref{fig:states}). A task is \emph{running} on the CPU, \emph{ready} to run and waiting its turn,
\emph{blocked} waiting for an event (a timer, a key, a disk block), or a \emph{zombie} that has finished but whose
resources have not been reclaimed yet.

\begin{figure}[htbp]
\centering
\begin{tikzpicture}[st/.style={draw=tkNavy!70, fill=tkLight, circle, minimum size=1.55cm, font=\sffamily\small\bfseries, align=center},
  every edge/.style={draw=tkNavy!80, -{Stealth[length=2mm]}, line width=0.7pt}, lab/.style={font=\sffamily\scriptsize, text=tkGray, align=center}]
  \node[st, fill=tkGray!10] (new) at (0,0) {new};
  \node[st, fill=tkGreen!12] (ready) at (3.4,0) {ready};
  \node[st, fill=tkRed!10] (run) at (7.4,0) {running};
  \node[st, fill=tkAmber!14] (blk) at (5.4,-2.8) {blocked};
  \node[st, fill=tkGray!18] (zom) at (10.6,0) {zombie};
  \draw (new) edge node[above, lab] {created} (ready);
  \draw (ready) edge[bend left=25] node[above, lab] {scheduled} (run);
  \draw (run) edge[bend left=25] node[below, lab] {time slice over\\or \code{yield()}} (ready);
  \draw (run) edge[bend left=15] node[right, lab] {\ sleep, wait for\\\ I/O or a lock} (blk);
  \draw (blk) edge[bend left=15] node[left, lab] {event happened\ \ } (ready);
  \draw (run) edge node[above, lab] {exit} (zom);
\end{tikzpicture}
\caption{Task states in Turk-OS. Only the scheduler moves a task into \emph{running}.}
\label{fig:states}
\end{figure}

\subsection{The context switch}

The key insight of \OSTEP chapter 6 is that the kernel regains control of the CPU in only two ways: the running code
calls into it (a system call, or \code{yield} for kernel threads), or a \textbf{timer interrupt} forces it in. Either
way, the kernel is now running on the current task's kernel stack. To switch, it saves the current task's registers,
records its stack pointer in its \code{struct task}, loads the next task's stack pointer, and restores that task's
registers. Because each task has its \textbf{own kernel stack}, the part of the old task's state that lives on the
stack (its return addresses, its interrupt frame) simply stays there until it is resumed.

The switch is done by a C-callable assembly function, \code{switch_context}. It only needs to save the four
callee-saved registers (EBX, ESI, EDI, EBP) plus ESP: under cdecl, the C code that calls it has already saved anything
else it cares about (Phase 1). When the new task's stack is loaded and \code{switch_context} executes \code{ret}, it
returns into wherever \emph{that} task called \code{switch_context} from, which might be the middle of a timer
interrupt from several seconds ago (Figure~\ref{fig:switch}).

\begin{figure}[htbp]
\centering
\begin{tikzpicture}[st/.style={draw=tkNavy!60, fill=tkLight, rounded corners=2pt, font=\sffamily\scriptsize, align=center, text width=4.4cm, minimum height=0.6cm},
  hi/.style={st, draw=tkRed!80!black, fill=tkRed!8}]
  \node[font=\sffamily\small\bfseries, text=tkNavy] at (0,0.7) {task A (being switched out)};
  \node[font=\sffamily\small\bfseries, text=tkNavy] at (7.2,0.7) {task B (being switched in)};
  \node[st] (a1) at (0,0) {A runs its code};
  \node[st, below=0.25cm of a1] (a2) {timer IRQ: CPU pushes EFLAGS, CS, EIP;\\\code{isr_common} saves all registers\\on A's kernel stack};
  \node[st, below=0.25cm of a2] (a3) {\code{timer_handler}: slice used up\\$\rightarrow$ \code{schedule()}};
  \node[hi, below=0.25cm of a3] (a4) {\code{switch_context(&A->esp, B->esp)}\\push EBP EBX ESI EDI; save ESP in A};
  \node[hi] (b1) at (7.2,-3.95) {load B's ESP; pop EDI ESI EBX EBP;\\\code{ret}};
  \node[st, above=0.25cm of b1] (b0) {B was switched out earlier inside its own \code{schedule()} call; its stack still holds that state};
  \node[st, below=0.25cm of b1] (b2) {back in B's \code{schedule()}, then B's\\\code{timer_handler}, \code{isr_common}};
  \node[st, below=0.25cm of b2] (b3) {\code{iret}: B continues exactly where it was interrupted};
  \draw[tkarrow] (a1) -- (a2); \draw[tkarrow] (a2) -- (a3); \draw[tkarrow] (a3) -- (a4);
  \draw[tkarrow, tkRed, line width=1pt] (a4.east) -- node[above, font=\sffamily\scriptsize, text=tkRed] {stack switch} (b1.west);
  \draw[tkarrow, dashed] (b0) -- (b1); \draw[tkarrow] (b1) -- (b2); \draw[tkarrow] (b2) -- (b3);
\end{tikzpicture}
\caption{A preemptive switch from A to B. A's half-finished interrupt stays frozen on A's stack until A is chosen again.}
\label{fig:switch}
\end{figure}

\subsection{Scheduling policy}

With the mechanism in place, the scheduler decides who runs. \OSTEP chapter 7 introduces the metrics:
\textbf{turnaround time} (from arrival to completion) and \textbf{response time} (from arrival to first run). Policies
like shortest-job-first optimise turnaround but need to know the future. \textbf{Round-robin} (RR) gives each ready
task a fixed \textbf{time slice} (quantum) in turn, which is excellent for response time and is what Turk-OS starts
with. Chapter 8's \textbf{multi-level feedback queue} (MLFQ) keeps several RR queues at different priorities, lowers
the priority of tasks that use their whole slice (CPU-bound work), and periodically boosts everyone to prevent
starvation. It favours interactive tasks without knowing anything about them in advance, which is why variants of it
ran in most UNIX systems. \OSDI\,\S2.4.5 adds the principle to keep policy and mechanism separate, so you can change
one without touching the other.

\section{Reading plan}

\begin{readingplan}
\must & \OSTEP & Ch.~4 \emph{The Abstraction: The Process} & 25--36 & Machine state, process states, the process list. \\
\must & \OSTEP & Ch.~6 \emph{Limited Direct Execution}, re-read \S6.3 \emph{Switching Between Processes} with its timeline & 52--58 & The exact mechanism of Tasks 8.3--8.6. \\
\must & \OSTEP & Ch.~7 \emph{Scheduling: Introduction}; Ch.~8 \emph{Scheduling: The Multi-Level Feedback Queue} & 63--86 & Metrics, RR, and the MLFQ rules for Task 8.9. \\
\must & \LOB & Ch.~14 \emph{Multitasking}: creating processes, cooperative and preemptive scheduling, separate kernel stacks & 73--75 & The x86 specifics, and the difficulties it warns about. \\
\should & \MOS & \S2.1 \emph{Processes} (model, creation, states, implementation) & 85--97 & Another view, with the interrupt-handling sequence spelled out. \\
\should & \MOS & \S2.4 \emph{Scheduling} & 148--167 & Goals per system type; RR, priority, multiple queues, lottery. \\
\should & \OSDI & \S2.5 \emph{Overview of Processes in MINIX 3}; \S2.6.5 \emph{Process Data Structures} & 112--125, 146--156 & MINIX's process table and scheduling queues. \\
\should & \OSDI & \S2.6.10 \emph{Scheduling in MINIX 3} & 182--185 & A real multi-queue scheduler in a few pages. \\
\should & \OSC & \S3.1--3.2 \emph{Process Concept}, \emph{Process Scheduling}; \S5.1--5.3 \emph{CPU scheduling} & 106--116, 200--217 & Queues, context switches, algorithms with examples. \\
\could & \OSTEP & Ch.~9 \emph{Proportional Share}; Ch.~10 \emph{Multiprocessor Scheduling} & 87--108 & Lottery and stride scheduling; what changes with more CPUs. \\
\could & \MOS & \S2.2 \emph{Threads}, especially \S2.2.5 \emph{Implementing Threads in the Kernel} & 97--113 & What you are building, placed among the alternatives. \\
\could & \OSC & \S5.7 \emph{Operating-System Examples} (Linux CFS, Windows) & 234--244 & How production schedulers do it. \\
\end{readingplan}

\begin{minixbox}
MINIX keeps each process's saved registers \emph{in the process table} (\code{p_reg} in \texttt{kernel/proc.h
(05500)}), not on a per-process kernel stack. Its assembly entry code in \texttt{kernel/mpx386.s (06300)} saves the
interrupted process's registers into its table entry, the kernel runs on one shared stack, and \code{restart} loads
whichever process was chosen. \code{enqueue}, \code{dequeue} and \code{pick_proc} in \texttt{kernel/proc.c (07400)}
manage 16 priority queues. This works for MINIX because its kernel never blocks in the middle of a call: work that has
to wait is done by user-space servers. Turk-OS gives every task its own kernel stack, like Linux and xv6, so a task can
block anywhere inside the kernel (in the middle of a disk read, for example) and resume there later.
\end{minixbox}

\section{Build it}

\task{The task structure}
Define \code{struct task} and a global list of all tasks (for \code{ps}), plus a \code{current} pointer. Keep the
saved stack pointer as the first field so the assembly never needs to know the struct layout. Hand out PIDs from a
counter.

\begin{lst}{c}{include/task.h}
enum task_state { TASK_READY, TASK_RUNNING, TASK_BLOCKED, TASK_ZOMBIE };

struct task {
    uint32_t         esp;          /* saved kernel ESP while not running (offset 0)  */
    uint32_t         pid;
    enum task_state  state;
    char             name[16];
    void            *kstack;       /* the stack allocation, freed when reaped        */
    uint32_t         wake_tick;    /* when a sleeping task should wake               */
    int              slice;        /* ticks left in the current time slice           */
    uint32_t         cpu_ticks;    /* accounting for ps                              */
    struct list_node node;         /* on the run queue, a sleep queue, or zombies    */
    struct list_node all;          /* on the list of all tasks                       */
};

extern struct task *current;
\end{lst}

\task{Kernel stacks with guard pages}
Give each task a 16\,KiB kernel stack in the region from \code{0xF0000000}, whose page tables already exist. Divide
the region into 20\,KiB slots and map only the top four pages of each slot, leaving the lowest page unmapped. A stack
overflow then hits the guard page and produces a page fault with an address just below the stack, instead of silently
overwriting the next task's memory. Keep freed slots on a list for reuse.
\verify{A test task that recurses without limit dies with a page fault whose address is in its own guard page.}

\task{switch\_context}
Type this in and comment it until you could explain every line to someone else.

\begin{lst}{asm}{kernel/arch/i386/switch.s}
; void switch_context(uint32_t *save_esp, uint32_t load_esp)
; Saves the current task's callee-saved registers and ESP, then resumes the task
; whose saved ESP is load_esp. Must be called with interrupts disabled.
global switch_context
switch_context:
    mov  eax, [esp + 4]     ; eax = &old->esp
    mov  edx, [esp + 8]     ; edx = new->esp
    push ebp                ; save callee-saved registers on the OLD task's stack
    push ebx
    push esi
    push edi
    mov  [eax], esp         ; old->esp = current stack pointer
    mov  esp, edx           ; switch to the NEW task's stack
    pop  edi                ; restore the new task's callee-saved registers
    pop  esi
    pop  ebx
    pop  ebp
    ret                     ; return to wherever the new task last called switch_context
\end{lst}

\task{Creating a task}
A new task has never called \code{switch_context}, so fake it: build its stack so the pops and \code{ret} land in a
small trampoline. The trampoline enables interrupts (the new task doesn't return through an interrupt path that
would restore them), calls the task's entry function, and calls \code{task_exit} if the function ever returns.

\begin{lst}{c}{kernel/proc/task.c}
extern void task_trampoline(void);          /* in switch.s, see below */

struct task *task_create(const char *name, void (*entry)(void *), void *arg)
{
    struct task *t = kcalloc(1, sizeof *t);
    t->kstack = kstack_alloc();             /* 16 KiB + guard page */
    uint32_t *sp = (uint32_t *)((char *)t->kstack + KSTACK_SIZE);

    *--sp = (uint32_t)task_trampoline;      /* switch_context's ret jumps here */
    *--sp = 0;                              /* ebp                             */
    *--sp = (uint32_t)entry;                /* ebx: the function to run        */
    *--sp = (uint32_t)arg;                  /* esi: its argument               */
    *--sp = 0;                              /* edi                             */
    t->esp = (uint32_t)sp;

    t->pid = next_pid++;
    strncpy(t->name, name, sizeof t->name - 1);
    make_ready(t);                          /* state = READY, append to run queue */
    return t;
}
\end{lst}

\begin{lst}{asm}{kernel/arch/i386/switch.s (continued)}
extern task_exit
global task_trampoline
task_trampoline:
    sti                     ; new tasks start with interrupts enabled
    push esi                ; argument
    call ebx                ; entry(arg)
    call task_exit          ; if entry returns, end the task (never returns)
\end{lst}
\verify{In GDB, break on \code{task_trampoline}: when it is hit, EBX holds your entry function and ESI its
argument.}

\task{Cooperative multitasking}
Turn the code that is already running (\code{kmain}) into task 0 by giving it a \code{struct task}; its
\code{esp} gets filled in at the first switch. Write \code{schedule()}, which moves the current task to the back of the
ready queue if it is still runnable and switches to the task at the front, and \code{yield()}, which calls it with
interrupts disabled. Save and restore the flags rather than blindly using \code{sti}, because \code{yield} may be
called with interrupts already off.

\begin{lst}{c}{kernel/proc/sched.c}
static struct list_node run_queue;          /* READY tasks, first in first out */

void schedule(void)                         /* interrupts must be disabled */
{
    struct task *prev = current;
    if (prev->state == TASK_RUNNING && prev != idle_task) {
        prev->state = TASK_READY;
        list_add_tail(&run_queue, &prev->node);
    }
    struct task *next = idle_task;
    if (!list_empty(&run_queue)) {
        next = container_of(run_queue.next, struct task, node);
        list_del(&next->node);
    }
    next->state = TASK_RUNNING;
    next->slice = QUANTUM_TICKS;            /* e.g. 5 ticks = 50 ms at 100 Hz */
    if (next == prev)
        return;
    current = next;
    switch_context(&prev->esp, next->esp);
}

void yield(void)
{
    uint32_t flags;
    __asm__ volatile ("pushf; pop %0; cli" : "=r"(flags));
    schedule();
    if (flags & 0x200)                      /* IF was set before: enable again */
        __asm__ volatile ("sti");
}
\end{lst}
\verify{Two tasks that each print a letter and call \code{yield()} in a loop produce \code{ABABAB...}; a third task
joins the rotation.}

\task{Preemption from the timer}
Now let the timer take the CPU away. In the timer handler, charge the tick to \code{current}, and when its slice
reaches zero, call \code{schedule()}. One change elsewhere is essential: in \code{isr_handler}, send the EOI to the PIC
\emph{before} calling an IRQ handler, not after. If the handler switches to another task, the old task's half-finished
interrupt stays frozen on its stack, and an EOI placed after the call would not be sent until that task runs again.
Until then the PIC delivers no more timer interrupts and preemption stops. Sending it first is safe because the
interrupt gate keeps interrupts disabled until \code{iret}.
\verify{Three tasks that loop printing without ever calling \code{yield} still take turns. The kernel monitor stays
responsive while they run.}

\task{Sleeping, idling and exiting}
Add \code{sleep_ms}: set \code{wake_tick}, mark the task blocked, put it on a sleep list, and call \code{schedule()}.
The timer handler moves every task whose \code{wake_tick} has passed back to the ready queue. Create an \code{idle}
task that runs \code{sti; hlt} forever; \code{schedule} picks it only when nothing else is ready, so the CPU halts
instead of spinning. Finally, \code{task_exit} marks the task as a zombie, puts it on a zombie list and schedules
away. It can't free its own stack, because it is still running on it, so another task (the idle task is a good
choice) reaps zombies: it frees their stacks and structures.
\verify{With every task sleeping, QEMU's CPU usage on the host drops to almost zero. Creating and exiting 1000 tasks
leaves the PMM and heap statistics exactly where they started.}

\task{ps and accounting}
Add a \code{ps} monitor command that walks the list of all tasks and prints PID, name, state and CPU ticks. Add
\code{spawn <name>} to start a demo task and \code{kill <pid>} to end one (for kernel threads, simply mark it to exit
at its next tick).

\task{A multi-level feedback queue (optional)}
Replace the single run queue with three queues and apply the rules of \OSTEP chapter 8: the highest non-empty queue
runs first; equal priorities share round-robin; new tasks start at the top; a task that uses up its allotment at a
level moves down; every second, all tasks move back to the top. Compare a CPU-bound task and an interactive one (the
monitor) with and without MLFQ.

\section{Testing and debugging}

When a switch goes wrong, the symptom usually appears later, in the new task. Break on \code{switch_context}, step
through it with \code{stepi}, and check with \code{x/6xw $esp} that the new stack holds exactly what the pops and
\code{ret} expect. The GDB command \code{info symbol <addr>} tells you where a return address points.

\begin{pitfalls}
Preemption works once, then the timer stops & EOI sent after the IRQ handler, which switched away & Send the EOI before calling the handler (Task~8.6). \\
New task crashes at its first instruction & The initial stack doesn't match \code{switch_context}'s pop order & Five words: EDI, ESI, EBX, EBP, return address, from low to high. \\
New task runs, but never gets preempted & Interrupts disabled in the new task & The trampoline must execute \code{sti}. \\
Random corruption, other tasks' data changes & Kernel stack overflow into a neighbour & Guard pages turn this into a page fault; make stacks bigger or recursion shallower. \\
Crash in \code{task_exit} or right after & The task freed its own stack while running on it & Only mark it as a zombie; let another task reap it. \\
Queue corruption under load & \code{schedule()} entered with interrupts enabled, so a timer interrupt re-entered it & Every path into \code{schedule()} disables interrupts first; assert IF is clear at its start. \\
\end{pitfalls}

\section{Milestone}

\begin{milestonebox}[Turk-OS multitasks]
\begin{checklist}
  \item Kernel threads run with their own 16\,KiB stacks and guard pages.
  \item \code{switch_context} and \code{task_trampoline} work; you can explain every instruction.
  \item At least three threads share the CPU preemptively, without calling \code{yield}.
  \item \code{sleep_ms} blocks properly; the idle task halts the CPU when nothing is ready.
  \item Tasks exit and are reaped; 1000 create/exit cycles leak no memory.
  \item \code{ps} shows every task with its state and CPU time.
\end{checklist}
\end{milestonebox}

\begin{questionbox}
\begin{enumerate}
  \item What state must be saved in a context switch, and why does \code{switch_context} save only four registers besides ESP?
  \item How does the timer interrupt make preemption possible, according to \OSTEP chapter 6?
  \item Draw the stack of a newly created task just before its first switch.
  \item Why must the EOI be sent before an IRQ handler that might call \code{schedule()}?
  \item Round-robin versus MLFQ: what problem does MLFQ solve, and why does it need a periodic priority boost?
  \item Why can't a task free its own kernel stack in \code{task_exit}?
  \item MINIX saves registers in the process table and uses a single kernel stack. What can Turk-OS do that MINIX's kernel can't, and what does it cost?
\end{enumerate}
\end{questionbox}

\section{Going further}

Implement lottery or stride scheduling from \OSTEP chapter 9 and compare fairness over a long run. Add a \code{top}
command that shows CPU usage per task over the last second. Measure the cost of a context switch by switching back and
forth a million times and reading the CPU's time-stamp counter (\code{rdtsc}) before and after. If you are heading for
the SMP track in Phase 15, read \OSTEP chapter 10 now and note what would have to change: one \code{current} per CPU,
per-CPU run queues, and real locks.

\nextphase{9}{Synchronization and Concurrency in the Kernel}
\booklegend
\end{document}
